\begin{abstract}\noindent\small
Reported property-type crime in Toronto changed in composition after 2019: auto theft rose while robbery fell.
We ask a narrow question: given that a major property-type crime was reported, which type is it?
Using \resG{n} Toronto Police incident records from 2014 to 2023, we contrast robbery with auto theft, break and enter, and theft over, and answer it twice.
A binomial logistic regression with neighbourhood-cluster-robust intervals describes how the odds of robbery vary by hour, month, weekday, solar darkness and a post-2019 period: the post-2019 odds ratio is \resG{post2019.or} \ci{\resG{post2019.robust_or_lo}}{\resG{post2019.robust_or_hi}}, and \resI{time_terms_clear} of \resI{time_terms} time terms have intervals that exclude one.
A five-rung prediction ladder, from time only to gradient-boosted trees, evaluated on held-out neighbourhoods and on a held-out later period, shows how little those time terms predict: under neighbourhood-grouped cross-validation time alone reaches a PR-AUC of \resA{variants.L1.pr_auc.estimate} \ci{\resA{variants.L1.pr_auc.lo}}{\resA{variants.L1.pr_auc.hi}} against a baseline of \resA{L0.pr_auc_baseline}, whereas adding place (premises, location type and coordinates) lifts it to \resA{variants.L2.pr_auc.estimate} \ci{\resA{variants.L2.pr_auc.lo}}{\resA{variants.L2.pr_auc.hi}} and census features add \resA{deltas.L3_minus_L2.pr_auc.estimate} \ci{\resA{deltas.L3_minus_L2.pr_auc.lo}}{\resA{deltas.L3_minus_L2.pr_auc.hi}}.
Statistically clear time associations therefore coexist with weak prediction from time alone.
Estimates describe the composition of reported incidents, not the risk of being robbed.
\end{abstract}
